% ==========================================================
% APÊNDICE — Lista de verificação SRQR
% ==========================================================

\section{Lista de Verificação SRQR (Standards for Reporting Qualitative Research)}
\label{app:srqr}

\noindent Lista de vinte e um itens proposta por \cite{OBrien2014} para o relato de pesquisa qualitativa, com indicação da seção do manuscrito em que cada item é atendido.

{\footnotesize
\setlength{\tabcolsep}{3pt}
\begin{longtable}{p{3.6cm}p{4.6cm}p{5.6cm}}
\toprule
\textbf{Item} & \textbf{O que a diretriz requer} & \textbf{Onde é atendido} \\
\midrule
\endfirsthead
\toprule
\textbf{Item} & \textbf{O que a diretriz requer} & \textbf{Onde é atendido} \\
\midrule
\endhead
\midrule
\multicolumn{3}{r}{\footnotesize\textit{continua na próxima página}}\\
\endfoot
\bottomrule
\endlastfoot
\textbf{1} Título & Identifica o estudo como qualitativo e indica abordagem, método ou tipo de dado & Título do capítulo, que nomeia o diagnóstico participativo e as categorias derivadas \\
\addlinespace
\textbf{2} Resumo & Sintetiza problema, questão, abordagem, método, resultados e conclusões & Resumo e Abstract \\
\addlinespace
\textbf{3} Formulação do problema & Descreve o problema, sua relevância e a lacuna na literatura & Introdução, com a lacuna enunciada em parágrafo próprio \\
\addlinespace
\textbf{4} Objetivo ou questão de pesquisa & Enuncia o objetivo, a questão ou a hipótese & Introdução, objetivo declarado em frase própria \\
\addlinespace
\textbf{5} Abordagem qualitativa e paradigma & Nomeia a abordagem e o paradigma que orientam o estudo & Delineamento Geral e Procedimento Analítico em Quatro Etapas, com Teoria Fundamentada em vertente construtivista \\
\addlinespace
\textbf{6} Características do pesquisador e reflexividade & Descreve quem conduziu o estudo, sua formação e sua relação com os participantes & Equipe de Pesquisa e Reflexividade, com composição disciplinar, via de acesso ao território e presença de terceiros nas entrevistas \\
\addlinespace
\textbf{7} Contexto & Caracteriza o local e as condições em que os dados foram produzidos & Área de Estudo, com localização, clima, base de recursos e critério de seleção das comunidades \\
\addlinespace
\textbf{8} Estratégia de amostragem & Explicita como e por que participantes e sítios foram selecionados & Área de Estudo, critérios de inclusão e exclusão; Entrevistas Semiestruturadas Abertas, critério de seleção dos entrevistados \\
\addlinespace
\textbf{9} Questões éticas & Informa aprovação ética, consentimento e proteção dos participantes & Aspectos Éticos, com CAAE, consentimento destacado para gravação, anuência coletiva em assembleia e tratamento conforme o Protocolo de Nagoia \\
\addlinespace
\textbf{10} Métodos de coleta & Descreve os tipos de dado, os procedimentos e a duração do trabalho de campo & Fase 1, entrevistas semiestruturadas abertas com duração declarada, oficinas de reconhecimento territorial \\
\addlinespace
\textbf{11} Instrumentos e tecnologias de coleta & Nomeia roteiros, equipamentos e programas utilizados & Entrevistas, com roteiro por eixos, transcrição automática no Transkriptor e codificação no ATLAS.ti \\
\addlinespace
\textbf{12} Unidades de estudo & Define o número e as características dos participantes, documentos e eventos & Fase 1, três entrevistas codificadas, oficinas em quatro comunidades e corpus documental de fonte prévia \\
\addlinespace
\textbf{13} Processamento dos dados & Descreve transcrição, saneamento, anonimização e gestão do material & Entrevistas, com transcrição integral, saneamento das transcrições automáticas, marcação de trechos inaudíveis e identificação por código \\
\addlinespace
\textbf{14} Análise dos dados & Explicita o processo analítico e quem o conduziu & Procedimento Analítico em Quatro Etapas e Guia de Códigos, e Confronto entre Categorias Emergentes e Campos do Questionário \\
\addlinespace
\textbf{15} Técnicas de confiabilidade & Nomeia as estratégias que sustentam a credibilidade dos achados & Parágrafo de critérios de rigor ao final da Fase 1, com triangulação, conferência pelo participante, planilha de rastreabilidade e cadeia de evidência \\
\addlinespace
\textbf{16} Síntese e interpretação & Apresenta os achados principais e sua interpretação & Resultados e Discussão, achados da Fase 1 \\
\addlinespace
\textbf{17} Vínculo com os dados empíricos & Sustenta cada achado em evidência do corpus & Citações diretas recuadas, tabela de categorias axiais e guia de códigos \\
\addlinespace
\textbf{18} Integração com trabalho prévio e contribuição & Relaciona os achados à literatura e discute transferibilidade & Triangulação Documental com Registros Locais Prévios e Especificações de Adaptação Produzidas na Fase 1 \\
\addlinespace
\textbf{19} Limitações & Discute as limitações de confiabilidade e alcance & Último parágrafo das Especificações de Adaptação, que delimita o alcance dos resultados \\
\addlinespace
\textbf{20} Conflitos de interesse & Declara conflitos potenciais e seu manejo & Declaração de conflito de interesses \\
\addlinespace
\textbf{21} Financiamento & Informa as fontes de financiamento e seu papel no estudo & Agradecimentos \\
\addlinespace
\end{longtable}
}
